\documentclass[12pt,a4paper]{article}
\usepackage[T2A]{fontenc}
\usepackage[utf8]{inputenc}
\usepackage[russian]{babel}
\usepackage{geometry}
\usepackage{tabularx} 
\geometry{margin=2.5cm}
\usepackage{amsmath,amssymb,mathtools}
\usepackage{enumitem}
\usepackage{tabularx}
\usepackage{hyperref}
\hypersetup{colorlinks=true,linkcolor=black,urlcolor=blue}
\setlength{\parindent}{0pt}
\setlength{\parskip}{6pt}
\begin{document}
{\Large\bfseries
\begin{flushleft}
Рубежный контроль №2 (Вариант №20)
\end{flushleft}}
\vspace{0.5em}\hrule\vspace{1.2em}

\section*{Задание №1}
\subsection*{Правило $bc\to cb$}
Это правило меняет местами соседние $b$ и $c$, сдвигая $c$ на одну позицию влево, и при этом это правило не меняет количества $b$ и $c$. Из слов вида $b^p c^q$, где все $b$ слева, а все $c$ справа), последовательными применениями этого правила можно получить любую перестановку букв $b$ и $c$

\subsection*{Правило $ac\to c^3$}
Это правило уменьшает количество букв a на 1 и увеличивает число c на 2.\\
В исходных базисных словах $a^n b^{n+k} c^k$ подстроки $ac$ нет. Чтобы применить $ac\to c^3$, надо сначала при помощи $bc\to cb$ протащить хотя бы один $c$ влево через все $b$, чтобы получить подстроку вида $a^n c\ldots$

По условию, базисным словом является $a^n b^{n+k}c^k$. Рассмотрю различные случаи при разных вариантах применении имеющихся правил.

\subsection*{Случай A: правило $ac\to c^3$ не применялось}
Меняется только порядок букв $b$ и $c$, при этом сохраняя их изначальное количество:
\[
|w|_b=n+k,,\qquad |w|_c=k
\]
Тогда могут получиться все слова вида
\[
L_1=\{\,a^n\,w \mid w\in\{b,c\}^*,\; |w|_b=n+k,\;|w|_c=k,\}.
\]

\subsection*{Случай B: правило $ac\to c^3$ применили $t\ge 1$ раз}
Это возможно только если $k\ge 1$, т.к. чтобы создать подстроку $ac$, нужен хотя бы один $c$.

Для применения этого правила можно протащить один $c$ сразу после $a^n$. В результате получим $a^n c\, b^{n+k} c^{k-1}$.

Теперь можно применить $ac\to c^3$ $t$ раз. Каждое применение уменьшает число $a$ на $1$ и добавляет $2$ к числу $c$.
После $t$ применений:
\[
|w|_a=n-t \quad 
\text{и появляется блок } c^{2t+1}.
\]
Оставшаяся часть слова состоит из букв $b$ и $c$ ($|w|_b=n+k,\quad |w|_c=k-1$) в любом порядке (их можно переставлять с помощью правила $bc\to cb$).

Итак,
\[
\begin{aligned}
L_2=\Bigl\{\,a^{n-t}\,c^{2t+1}\,w \;\Bigm|\;& n\ge 0,\;k\ge 1,\;1\le t\le n,\\
& w\in\{b,c\}^*,\; |w|_b=n+k,\;|w|_c=k-1 \Bigr\}.
\end{aligned}
\]
Отсюда следует, что 
\[
L = L_1 \cup L_2.
\]

\subsection*{Представление языка в виде PDA}
\begin{center}
\begin{minipage}{\linewidth}
\centering
\includegraphics[width=1.0\linewidth]{PDA.jpg}
\captionof{Рис. 1: PDA-автомат}
\label{fig:mindfa.}
\end{minipage}
\end{center}

\subsection*{Проверка регулярности языка }
Рассмотрим регулярный язык
\[
R = a^*b^*,
\]
Пересеку его с имеющимся языком L
\[
L \cap R = \{a^n b^n \mid n\ge 0\}.
\]

Язык
\[
\{a^n b^n \mid n\ge 0\}
\]
не является регулярным. Это можно легко доказать в том числе используя таблицу классов эквивалентности.
Если бы $L$ был регулярным, то и пересечение $L\cap R$ было бы регулярным. Следовательно, язык L не является регулярным

\newpage

\section*{Задание №2}

\subsection*{Рассмотрим случаи, в которых точно можно быть уверенным, что слово будет принадлежать языку}

1. Пусть \( |w_1| = 2 \) и \( w_1^R \) найдено больше, чем 1 раз. Тогда из этих \( (w_1)^R \) можно выбрать подходящий вариант, чтобы было выполнено условие \( |z_1| \neq |z_2| \).

2. Пусть \( |w_1| = 3 \) и найдено \( w_1^R \). Тогда при необходимости вместо \( |w_1| = 3 \) можно выбрать \( |w_1| = 2 \), в результате чего точно будет выполняться условие \( |z_1| \neq |z_2| \).

\subsection*{Случай, где пункты 1 и 2 не выполняются}

Если пункты 1 и 2 не выполняются, то тогда известно, что \( |w_1| = 2 \) и \( w_1^R \) единственно. Для решения этого случая можно нарисовать DPDA.

Рассмотрим все возможные тройки, с которых начинается слово (aaa, aab, aba, abb, baa, bab, bba, bbb):

\begin{center}
\begin{minipage}{\linewidth}
\centering
\includegraphics[width=1.0\linewidth]{startdpda.jpg}
\captionof{Рис. 2: Начало общего DPDA}
\label{fig:startdpda.jpg}
\end{minipage}
\end{center}

Каждая из этих вершин будет иметь уникальный DPDA для поиска развёрнутого пути до этой вершины. Так, одна из веток (для тройки \( aba \)) будет выглядеть следующим образом:

\begin{center}
\begin{minipage}{\linewidth}
\centering
\includegraphics[width=1.0\linewidth]{branchdpda.jpg}
\captionof{Рис. 3: ветвь DPDA по пути aba}
\label{fig:branchdpda.jpg}
\end{minipage}
\end{center}

Используя те же идеи, можно построить DPDA для каждой тройки. В результате этого получится общий DPDA для этого языка.\\
Таким образом можно утверждать, что язык является DCFL.

\subsection*{Проверка LL-свойства}
Рассмотрим слово, входящее в исходный язык: 
\(abab^{(n+k)}aa\) .
Для \(n\) подбирается такое значение, что после чтения \(n\) символов стек будет содержать не менее, чем \(k+3\) стековых символа. Рассмотри еще одно слово, входящее в исходный язык: 
\(abab^{(n+k)}aab^{(n+k+1)}\).
Для данного слова первое слово \(abab^{(n+k)}aa\) является префиксом, к которому был добавлен суффикс \(b^{(n+k+1)}\).
При чтении фрагмента \(b^{k}aa\) автомат может снять со стека не более \(k+2\) символов. Тогда при чтении слова  \(abab^{(n+k)}aa\) стек гарантированно будет непустым, и автомат не сможет корректно завершить работу. Получается неожиданный конец при том, что префикс \(abab^{(n+k)}aa\) продолжим до слова из \(L\). Следовательно, язык не
может быть \(LL(k)\).

\newpage

\section*{Задание №3}
Из первых двух правил для $S$:
\[
S \to TbS,\qquad S \to T
\]
следует, что любой вывод из $S$ имеет вид
\begin{align*}
S &\Rightarrow TbS\\
  &\Rightarrow TbTbS\\
  &\Rightarrow \cdots\\
  &\Rightarrow (Tb)^nT.
\end{align*}

То есть любое слово языка представимо в форме
\[
w = u_0\,b\,u_1\,b\,u_2\,b\,\cdots\,b\,u_n,
\qquad \text{где } u_i \in L(T).
\]
Важно отметить, что буквы $b$ между блоками это именно те $b$, которые получаются из правила $S \to TbS$.

Используем семантику правила
\[
S \to T\,b\,S
\quad;\quad
S_0.a := S_1.a + 1,\qquad S_1.a == T.a.
\]
то есть:
\begin{itemize}
  \item атрибут хвоста $S_1.a$ должен совпасть с атрибутом текущего $T.a$;
  \item атрибут ``всего'' на $1$ больше атрибута хвоста.
\end{itemize}

Если записать слово как $u_kbu_{k-1}b\ldots bu_0$, то получается:
\begin{itemize}
  \item правый блок $u_0$ (последний $T$) получает некоторое значение
  \[
  T_0.a = f;
  \]
  \item блок перед ним обязан иметь то же самое значение (потому что хвост равен текущему $T.a$):
  \[
  T_1.a = f;
  \]
  \item следующий левее уже должен быть на 1 больше:
  \[
  T_2.a = f + 1;
  \]
  \item дальше по цепочке каждый левый блок на 1 больше предыдущего:
  \[
  T_3.a = f + 2,\ \ldots,\ T_k.a = f + (k - 1).
  \]
\end{itemize}

\subsection*{Правило $T\to aTa$}
Единственное правило для $T$, которое порождает $a$, имеет вид
\[
T \to aTa.
\]
Следовательно, каждое применение правила $T\to aTa$ добавляет ровно две буквы $a$. Поэтому любое слово, выводимое из $T$ и состоящее только из букв $a$, имеет вид
\[
a^{2m}
\]
для некоторого $m\ge 0$. То есть степень всегда кратна $2$.

Чтобы вывести слово $a^{2m}$ из $T$, правило $T\to aTa$ необходимо применить ровно $m$ раз. После чего вывод заканчивается правилом $T\to \varepsilon$.

Для правила $T\to aTa$ возможны две семантики:
\[
T_0.a := T_1.a + 1
\qquad \text{или} \qquad
T_0.a := T_1.a + 2.
\]

Значит, итоговое значение $T.a$ на слове $a^{2m}$ равно сумме $m$ слагаемых, каждое из которых равно $1$ или $2$. Отсюда:
\[
\min T.a = m \quad (\text{если все $m$ раз брать $+1$}), \qquad
\max T.a = 2m \quad (\text{если все $m$ раз брать $+2$}),
\]
таким образом
\[
T.a \in [m,\,2m].
\]

\subsection*{Выбор слова для накачки}
Рассмотрим слово
\[
w_n = a^{2n+2}\,b\,a^{4n}\,b\,a^{2n}.
\]
Оно имеет вид $T_2\,b\,T_1\,b\,T_0$, где
\[
T_0 \Rightarrow a^{2n},\qquad
T_1 \Rightarrow a^{4n},\qquad
T_2 \Rightarrow a^{2n+2}.
\]
Следовательно, число применений правила $T\to aTa$ равно:
\[
\text{для }T_0:\ n\ \text{раз},\qquad
\text{для }T_1:\ 2n\ \text{раз},\qquad
\text{для }T_2:\ n+1\ \text{раз}.
\]

Чтобы слово входило в язык, должны существовать такие значения, что
\[
T_0.a=f,\qquad T_1.a=f,\qquad T_2.a=f+1.
\]

Так как $T_0\Rightarrow a^{2n}$, то по описанному ранее имеем, что
\[
T_0.a \in [n,\,2n].
\]
Так как $T_1\Rightarrow a^{4n}=a^{2\cdot(2n)}$, то по описанному ранее имеем, что
\[
T_1.a \in [2n,\,4n].
\]
$T_0.a=T_1.a=f$, поэтому
\[
f\in [n,\,2n]\cap[2n,\,4n]=\{2n\},
\]
Таким образом происходит фиксация значения для $T_0.a$ и $T_1.a$

$T_2 \Rightarrow a^{2n+2}=a^{2\cdot(n+1)}$, тогда
\[
T_2.a \in [n+1,\,2n+2].
\]
Поскольку
\[
2n+1 \in [n+1,\,2n+2],
\]
то нужное значение для $T_2.a$ действительно допустимо, и при этом сохраняется пересечение с пересечением $T_1.a$ и $T_0.a$

Необходимо добиться
\[
T_0.a=2n,\qquad T_1.a=2n,\qquad T_2.a=2n+1.
\]
Для $T_0\Rightarrow a^{2n}$.
Правило $T\to aTa$ применяется $n$ раз.
Чтобы получить $T_0.a=2n$, на каждом из $n$ применений берём семантику $+2$:\\
Для $T_1\Rightarrow a^{4n}$.
Правило $T\to aTa$ применяется $2n$ раз.
Чтобы получить $T_1.a=2n$, на каждом из $2n$ применений берём семантику $+1$:\\
Для $T_2\Rightarrow a^{2n+2}$.
Правило $T\to aTa$ применяется $n+1$ раз.
Чтобы получить $T_2.a=2n+1$, берём $n$ раз семантику $+2$ и один раз семантику $+1$:\\

Следовательно,
\[
w_n = a^{2n+2}\,b\,a^{4n}\,b\,a^{2n}\in L.
\]

\subsection*{Накачка}

\begin{enumerate}
  \item \textbf{Отрицательная накачка внутри \(a^{2n+2}\)}\\
  Первый блок становится короче: \(a^{2n+2}\mapsto a^{2n+2-k}\).
  Для такого блока возможные значения атрибута лежат в диапазоне
  \(\;T_2.a\in[n+1-k,\;2(n+1-k)]\), то есть максимум \(\le 2n\).
  Т.к. для этого слова обязательно должно быть \(\;T_2.a = 2n+1\)- слово выходит из языка.

  \item \textbf{Положительная накачка внутри \(a^{4n}\)}\\
  Второй блок становится длиннее: \(a^{4n}\mapsto a^{4n+k}\).
  Тогда минимально возможное значение становится \(\;T_1.a \ge 2n+k > 2n\).
  А правый блок \(a^{2n}\) даёт значения только из \([n,2n]\), поэтому
  равенство \(\;T_1.a = T_0.a\) выполнить нельзя — слово выходит из языка.

  \item \textbf{Отрицательная накачка внутри \(a^{2n}\)}\\
  Третий блок становится короче: \(a^{2n}\mapsto a^{2n-k}\).
  Тогда максимально возможное значение \(\;T_0.a \le 2n-k < 2n\),
  а средний блок \(a^{4n}\) даёт значения \(\;T_1.a \in [2n,4n]\), то есть \(\;T_1.a\ge 2n\).
  Следовательно, равенство \(\;T_1.a = T_0.a\) невозможно — слово выходит из языка.

  \item \textbf{Накачка попарно внутри $a^{2n+2}$ и $a^{4n}$}\\
  Положительная накачка блока $a^{4n}$ увеличивает минимальное и максимальное значения, и тогда снова $T_1.a$ становится больше $2n$, а $T_0.a$ остаётся $\le 2n$ - слово выходит из языка.

  \item \textbf{Накачка попарно внутри $a^{4n}$ и $a^{2n}$}\\
  Положительная накачка блока $a^{2n}$ при больших значениях приведет к тому, что минимальное значение $T_0.a$ будет больше минимального значения для $T_1.a$, в результате чего между ними пропадёт пересечение - слово выходит из языка
\end{enumerate}

\subsection*{Противоречия по $b$ (из правила $T \to bT$)} 
Нужно проверить, может ли какая-то из букв $b$ в этом слове появиться за счёт правила $T \to bT$.
Если буква $b$ появилась именно из правила $T \to bT$ и перед этой $b$ в итоговом слове стоят только $a$, то эти $a$ могли появиться только из обёрток $T \to aTa$, а значит ровно столько же $a$ обязано стоять и в конце того же $T$-блока.

\textbf{Первая $b$ (между $a^{2n+2}$ и $a^{4n}$) получена правилом}\\
Чтобы попытаться получить слово $w$ при укзаанных условиях единственный порядок действий выглядит следующим образом:
\[
\begin{aligned}
& S \to TbS \to TbT \to Tba^nTa^n \to Tba^{2n} \to a^{2n+2}Ta^{2n+2}ba^{2n} \to \\
& \to a^{2n+2}bTa^{2n+2}ba^{2n} \to a^{2n+2}ba^{n-1}Ta^{n-1}a^{2n+2}ba^{2n} \to a^{2n+2}ba^{4n}ba^{2n}
\end{aligned}
\]
при выполнени такого порядка действий промежуток значений $T_0.a = [n,\, 2n]$, тогда как промежуток значений $T_1.a = [2n+2+2\cdot(n-1),\, 4n+4+2\cdot(2n-2)] = [4n,\, 8n]$. Значения блоков должны быть равны, а поскольку промежутки не имеют пересечения, это слово в языке таким образом не могло быть получено

\textbf{Вторая $b$ (между $a^{4n}$ и $a^{2n}$) получена правилом}\\
Чтобы попытаться получить слово $w$ при укзаанных условиях единственный порядок действий выглядит следующим образом:
\[
\begin{aligned}
& S \to TbS \to TbT \to a^{n+1}Ta^{n+1}bT \to a^{2n+2}ba^{4n}Ta^{4n} \\
& \to a^{2n+2}ba^{4n}bTa^{4n} \to a^{2n+2}ba^{4n}ba^{4n}
\end{aligned}
\]
$4n>2n$, следовательно это слово в языке таким образом не могло быть получено

\textbf{Обе $b$ получены правилом}\\
Чтобы попытаться получить слово $w$ при укзаанных условиях единственный порядок действий выглядит следующим образом:
\[
\begin{aligned}
& S \to T \to a^{2n+2}Ta^{2n+2} \to a^{2n+2}bTa^{2n+2} \to a^{2n+2}ba^{4n}Ta^{4n}a^{2n+2} \\
& \to a^{2n+2}ba^{4n}bTa^{6n+2} \to a^{2n+2}ba^{4n}ba^{6n+2}
\end{aligned}
\]
$6n+2>2n$, следовательно это слово в языке таким образом не могло быть получено\\
\\
Аналогичные результаты получаются при исследовании противоречий для каждой накачки.\\
\\
Таким образом можно сделать вывод, что данный язык не является КС-языком

\end{document}
